\section{Proofs}
\label{appendix:proofs}

This appendix gives the formal support for the four claims summarized in
\S\ref{sec:method-theory}: reward boundedness (Prop.~\ref{prop:reward-bounded}),
posterior consistency in the stationary no-discount regime
(Prop.~\ref{prop:consistency}), a stationary well-specified regret reduction
(Thm.~\ref{thm:stationary-regret}), and a fixed-schedule adaptation rate after
an abrupt change (Thm.~\ref{thm:adaptation}). The regret theorem is intentionally
stated for a clean version of the score; \S\ref{app:proof-limitations} records
the gap to the full empirical setting.

\subsection{Notation and assumptions (recap)}
\label{app:proof-setup}

We use the notation of \S\ref{sec:problem} and \S\ref{sec:method}: $K$ arms,
$K_b$ context buckets, round $t \in \{0,\ldots,T-1\}$, bucket
$k_t=k(\mathbf{x}_t)$, observed quality $q_t \in [0,1]$, failure flag
$f_t\in\{0,1\}$, realized cost $c_t\geq0$, utility
$\ut=q_t-\nu f_t$ with $\nu=0.5$, normalized cost
$\ctilde=\mathrm{clip}(c_t/c_{\max},0,1)$, wallet-pressure weight
$\lamn=\lambda_t/(1+\lambda_t)\in(0,1)$, and payment-aware reward
$\rt=(1-\lamn)\ut-\lamn\ctilde$.
All regret sums below are over served rounds with nonempty affordable set; if a
wallet exhausts early, replace $T$ by the served horizon $\tau$ of
\S\ref{sec:problem}.

For each cell $(a,k)$, PA-DCT keeps discounted sufficient statistics
$(n^{a,k}_t,S^{a,k}_t,{S_c}^{a,k}_t)$ and Normal--Normal Q/C posteriors.
It is useful to write the posterior means as
\begin{equation}
\label{eq:posterior-mean-compact}
\thetaq{a}{k}
= \frac{\kappa_q\mu_{0,q}+S^{a,k}_t}{\kappa_q+n^{a,k}_t},
\qquad
\thetac{a}{k}
= \frac{\kappa_c\bar c_a+{S_c}^{a,k}_t}{\kappa_c+n^{a,k}_t},
\end{equation}
where $\kappa_q=\sigma_q^2/\sigma_{0,q}^2$ and
$\kappa_c=\sigma_c^2/\sigma_{0,c}^2$.
We write a single $n^{a,k}_t$ and $\gamma$ because the locked experiments use
$\gamma_c=\gamma$ and each selected call supplies one utility and one cost
observation. If quality and cost discounts are split, the same formulas apply to
the two posteriors with their own effective counts.

Throughout this appendix we use:
\begin{description}
\item[(A1)] \emph{Stationary or piecewise-stationary cells.} Conditional on the
chosen arm $a$ and bucket $k$, and averaging over the exogenous task draw within
that bucket, the joint distribution of $(q_t,c_t,f_t)$ is identical within each
phase.
\item[(A2)] \emph{Sub-Gaussian noise.} The utility $\ut$ given cell $(a,k)$ is
$\sigma_u$-sub-Gaussian around mean $\mu_u^{a,k}$; the realized cost $c_t$ is
$\sigma_{c,\mathrm{true}}$-sub-Gaussian around mean $\mu_c^{a,k}$. The policy
uses fixed proxy variances $\sigma_q^2,\sigma_c^2$, possibly different from the
true variances.
\item[(A3)] \emph{Bounded means.} $|\mu_u^{a,k}|\leq1$ and
$\mu_c^{a,k}\in[0,c_{\max}]$ for every $(a,k)$.
\end{description}

The stationary regret reduction further uses:
\begin{description}
\item[(A2$^\prime$)] \emph{Well-specified product model.} The policy's
proxy likelihoods are the true likelihoods: utility and cost observations are
conditionally independent Gaussians around $(\mu_u^{a,k},\mu_c^{a,k})$ with
variances $\sigma_q^2=\sigma_u^2$ and
$\sigma_c^2=\sigma_{c,\mathrm{true}}^2$, and the prior factorizes as the
policy's independent Gaussian Q/C priors.
\item[(A3$^\prime$)] \emph{Clean score.} The regret reduction is stated for the
unclipped score
$(1-\lamn)\hatu{a}{k}-\lamn\hatc{a}{k}/c_{\max}$, or equivalently for
trajectories on which the implementation's cost-sample clip in
Eq.~\eqref{eq:padct-score} is inactive.
\end{description}
Propositions~\ref{prop:reward-bounded}--\ref{prop:consistency} and
Theorem~\ref{thm:adaptation} use only (A1)--(A3); only
Theorem~\ref{thm:stationary-regret} uses (A2$^\prime$)--(A3$^\prime$).

\subsection{Proposition~\ref{prop:reward-bounded}: reward boundedness}
\label{app:proof-bounded}

\begin{proposition}[Reward boundedness]
\label{prop:reward-bounded}
For every round $t$, the payment-aware reward satisfies $\rt\in[-1,+1]$.
\end{proposition}

\begin{proof}
Since $q_t\in[0,1]$, $f_t\in\{0,1\}$, and $\nu=0.5$, we have
$|\ut|=|q_t-0.5f_t|\leq1$; also $\ctilde\in[0,1]$ and
$\lamn\in(0,1)$. Hence
$|\rt|\leq(1-\lamn)|\ut|+\lamn|\ctilde|\leq(1-\lamn)+\lamn=1$.
\end{proof}

\subsection{Proposition~\ref{prop:consistency}: posterior consistency}
\label{app:proof-consistency}

\begin{proposition}[Posterior consistency, stationary regime]
\label{prop:consistency}
Fix $(a,k)$ and assume the utility distribution at $(a,k)$ is stationary with
finite variance and mean $\mu_u^{a,k}$. With no discount ($\gamma=1$), the
Q-posterior mean $\thetaq{a}{k}$ converges almost surely to $\mu_u^{a,k}$ as the
cell's pull count $n^{a,k}\to\infty$. The same holds for the C-posterior and
$\mu_c^{a,k}$.
\end{proposition}

\begin{proof}
Let $\{\ut^{(i)}\}_{i=1}^n$ be the utility observations in cell $(a,k)$ after
$n$ pulls. By (A1) and the strong law of large numbers,
$S^{a,k}_t/n^{a,k}_t=n^{-1}\sum_i\ut^{(i)}\to\mu_u^{a,k}$ almost surely.
Dividing the first expression in~(\ref{eq:posterior-mean-compact}) by
$n^{a,k}_t$ in numerator and denominator gives
$\thetaq{a}{k}\to\mu_u^{a,k}$ because $\kappa_q/n^{a,k}_t\to0$. The cost
posterior is identical with $\bar c_a,{S_c}^{a,k}_t,\kappa_c$ in place of
$\mu_{0,q},S^{a,k}_t,\kappa_q$.
\end{proof}

\subsection{Lemma: discounted effective sample size}
\label{app:lemma-ess}

\begin{lemma}[Effective sample size]
\label{lem:ess}
For any $\gamma\in(0,1)$ and any sequence of pulls of cell $(a,k)$,
\[
n^{a,k}_t \leq \frac{1-\gamma^{t+1}}{1-\gamma}\leq\frac{1}{1-\gamma}.
\]
The first bound is tight when $(a,k)$ is pulled at every round.
\end{lemma}

\begin{proof}
If $\mathcal{T}^{a,k}_t\subseteq\{0,\ldots,t\}$ is the set of pull rounds for
cell $(a,k)$, the round-based discount rule gives
$n^{a,k}_t=\sum_{s\in\mathcal{T}^{a,k}_t}\gamma^{t-s}
\leq\sum_{j=0}^t\gamma^j=(1-\gamma^{t+1})/(1-\gamma)$.
\end{proof}

\subsection{Theorem~\ref{thm:stationary-regret}: stationary regret reduction}
\label{app:proof-stationary}

\begin{theorem}[Clean-score regret reduction to Gaussian Thompson sampling]
\label{thm:stationary-regret}
Under (A1), (A2$^\prime$), (A3$^\prime$), and $\gamma=1$, define the
\emph{policy-state PA-regret} of clean-score PA-DCT against the wallet-coupled
oracle as
\begin{equation}
\label{eq:policy-state-regret}
\begin{aligned}
g_t(a) &:=
(1-\lamn)\,\mu_u^{a,k_t}
- \lamn\,\mu_c^{a,k_t}/c_{\max},\\
\Reg_T &:= \sum_{t=0}^{T-1}\big(g_t(a^\star_t)-g_t(a_t)\big),
\end{aligned}
\end{equation}
where $a^\star_t\in\arg\max_{a\in\mathcal{A}_t}g_t(a)$ at the policy's wallet
pressure $\lamn$, and $a_t$ is PA-DCT's chosen arm. Conditional on the
bucket-visit sequence, clean-score PA-DCT reduces to per-bucket Gaussian
Thompson sampling with a time-varying linear context. For a bucket visited
$T_k$ times,
\begin{equation}
\label{eq:cell-regret}
\mathbb{E}\!\left[\Reg^{\mathrm{cell}\,k}_{T_k}\right]
\leq C_0\sqrt{K\,T_k\,\log T_k},
\end{equation}
for a constant $C_0$ depending on the prior and likelihood variances and on
$c_{\max}$. Aggregating across $K_b$ buckets with $\sum_kT_k=T$ gives
\begin{equation}
\label{eq:stationary-regret}
\mathbb{E}\!\left[\Reg_T\right]\leq C\sqrt{K\,K_b\,T\log T}.
\end{equation}
\end{theorem}

\begin{proof}
For a visit to bucket $k$, write $w_t=\lamn$ and define the observed feature
\[
\phi_t=(1-w_t,\,-w_t/c_{\max}),\qquad
\theta^{a,k}=(\mu_u^{a,k},\,\mu_c^{a,k}).
\]
Then $g_t(a)=\phi_t^\top\theta^{a,k}$. Under (A2$^\prime$), the independent
Q/C Thompson samples used by PA-DCT imply
\[
\begin{aligned}
&(1-w_t)\hatu{a}{k}-w_t\hatc{a}{k}/c_{\max}\\
&\quad\sim
\mathcal{N}\!\left(
\phi_t^\top(\thetaq{a}{k},\thetac{a}{k}),
(1-w_t)^2\Sigmaq{a}{k}
+w_t^2\Sigmac{a}{k}/c_{\max}^2
\right),
\end{aligned}
\]
which is exactly the product-Gaussian posterior sample of
$\phi_t^\top\theta^{a,k}$. Thus the clean-score policy is Gaussian Thompson
sampling on a linear payoff with a two-dimensional, time-varying feature.

The feature $\phi_t$, bucket $k_t$, and affordable set $\mathcal{A}_t$ are
observed before sampling. Bucket visits are determined by the exogenous task
stream in \S\ref{sec:problem}, so conditioning on a bucket subsequence does not
introduce policy-dependent context selection; the wallet pressure and affordable
set are predictable state shared by the policy-state comparator in
(\ref{eq:policy-state-regret}). Therefore the per-round gap is the standard
contextual Thompson-sampling regret for at most $K$ arms.

Applying the Bayesian regret bound for Gaussian-prior Thompson sampling
\citep[Theorem~1]{russo2014learning} to each bucket subsequence (equivalently,
the K-armed Gaussian-prior specialization in \citealt[Ch.~36]{lattimore2020bandit})
gives~(\ref{eq:cell-regret}). Summing over buckets and using Cauchy--Schwarz,
$\sum_k\sqrt{T_k\log T_k}\leq\sqrt{K_bT\log T}$, gives
(\ref{eq:stationary-regret}).
\end{proof}

\subsection{Theorem~\ref{thm:adaptation}: adaptation after an abrupt change}
\label{app:proof-adaptation}

\begin{theorem}[Adaptation rate under an abrupt change, fixed schedule]
\label{thm:adaptation}
Fix $\gamma\in(0,1)$ and cell $(a,k)$. Suppose the true utility mean shifts
between rounds $t^\star-1$ and $t^\star$, with pre-shock mean
$\mu_{\mathrm{old}}$, post-shock mean $\mu_{\mathrm{new}}$, and
$\Delta=|\mu_{\mathrm{old}}-\mu_{\mathrm{new}}|$. Consider any fixed, or
reward-noise-independent, post-shock pull schedule
$\{s_1<\cdots<s_m\}\subseteq[t^\star,t]$ for this cell. Define
\[
A:=\gamma^{t-t^\star+1}n^{a,k}_{t^\star-1},\qquad
B:=\sum_{i=1}^m\gamma^{t-s_i},\qquad
\kappa_q:=\sigma_q^2/\sigma_{0,q}^2 .
\]
Then the Bayesian Q-posterior mean satisfies
\begin{equation}
\label{eq:adaptation}
\left|\,\mathbb{E}[\theta^{a,k}_{q,t}]-\mu_{\mathrm{new}}\,\right|
\leq
\frac{A\Delta+\kappa_q|\mu_{0,q}-\mu_{\mathrm{new}}|}
{\kappa_q+A+B},
\end{equation}
where the expectation is over observation noise for the fixed schedule.
\end{theorem}

\begin{remark}[Two regimes]
\label{rem:adapt-regimes}
The bound~(\ref{eq:adaptation}) separates a prior term from a stale-data term.
When $\kappa_q\gg A+B$, the prior controls the bias. When
$\kappa_q\ll A+B$, the leading term is $\Delta A/(A+B)$, which is reduced by
discounting and by new post-shock pulls.
\end{remark}

\begin{corollary}[Recovery horizon, data-dominated regime]
\label{cor:recovery}
Suppose $(a,k)$ is pulled at every post-shock round, so
$\tau:=t-t^\star+1=m$, and assume the data-dominated regime
$\kappa_q\ll A+B$. Then
\[
\left|\,\mathbb{E}[\theta^{a,k}_{q,t}]-\mu_{\mathrm{new}}\,\right|
\lesssim
\Delta\gamma^\tau+O\!\big(\kappa_q/(A+B)\big),
\]
and the leading stale-data term is at most $\delta$ once
\[
\tau\geq\frac{\log(\Delta/\delta)}{\log(1/\gamma)}
=O\!\big(\log(\Delta/\delta)/(1-\gamma)\big).
\]
\end{corollary}

\begin{proof}[Proof of Theorem~\ref{thm:adaptation}]
For the fixed schedule, conditioning on the pull times does not change the mean
of the included observations. From~(\ref{eq:posterior-mean-compact}),
\[
\theta^{a,k}_{q,t}=\frac{\kappa_q\mu_{0,q}+S^{a,k}_t}
{\kappa_q+n^{a,k}_t}.
\]
Splitting the discounted mass into pre- and post-shock pulls gives
$n^{a,k}_t=A+B$ and
$\mathbb{E}[S^{a,k}_t]=A\mu_{\mathrm{old}}+B\mu_{\mathrm{new}}$. Hence
\[
\mathbb{E}[\theta^{a,k}_{q,t}]-\mu_{\mathrm{new}}
=
\frac{\kappa_q(\mu_{0,q}-\mu_{\mathrm{new}})
+A(\mu_{\mathrm{old}}-\mu_{\mathrm{new}})}
{\kappa_q+A+B}.
\]
Taking absolute values and using
$|\mu_{\mathrm{old}}-\mu_{\mathrm{new}}|=\Delta$ gives
(\ref{eq:adaptation}).
\end{proof}

\begin{proof}[Proof of Corollary~\ref{cor:recovery}]
If the cell is pulled at every post-shock round, then
$B=\sum_{j=0}^{\tau-1}\gamma^j=(1-\gamma^\tau)/(1-\gamma)$ and
$A=\gamma^\tau n^{a,k}_{t^\star-1}$. By Lemma~\ref{lem:ess},
$n^{a,k}_{t^\star-1}\leq1/(1-\gamma)$, so
\[
\frac{A}{A+B}
\leq
\frac{\gamma^\tau/(1-\gamma)}
{\gamma^\tau/(1-\gamma)+(1-\gamma^\tau)/(1-\gamma)}
=\gamma^\tau .
\]
Substituting this into~(\ref{eq:adaptation}) yields the displayed bound; solving
$\Delta\gamma^\tau\leq\delta$ gives the stated horizon.
\end{proof}

\begin{remark}[Cost-posterior analogue]
The same calculation applies to realized-cost shocks by replacing
$(\mu_{\mathrm{old}},\mu_{\mathrm{new}},\mu_{0,q},\kappa_q,\gamma)$ with
$(\mu^c_{\mathrm{old}},\mu^c_{\mathrm{new}},\bar c_a,\kappa_c,\gamma_c)$ and
using the C-posterior statistic. We state the utility version to avoid
duplicating notation; this is the algebra used for S3-style price changes once
the changed arm is pulled again.
\end{remark}

\subsection{Scope of the formal analysis}
\label{app:proof-limitations}

The results above are intended as component-level formal support rather than a
complete regret theorem for the full replay system. The stationary regret
reduction assumes a well-specified product-Gaussian model, $\gamma=1$, and the
clean score with inactive cost-sample clipping; the implemented clipped scorer
coincides with this policy on trajectories where the clip is inactive. The
adaptation result is conditional on a fixed post-shock pull schedule, so it
explains how discounted posteriors forget stale evidence after a changed arm is
revisited, but does not prove an unconditional Thompson-sampling revisit rate or
a cumulative piecewise-stationary regret bound. Finally, the theorem's comparator
uses the policy's own wallet pressure and affordable set, whereas the empirical
\emph{True Oracle} has its own wallet trajectory.
